\documentclass[../../main.tex]{subfiles}
\graphicspath{{../../image/}{../../figure/}}
\begin{document}

\section{弹性动力学中的空化}\label{section:ch18-弹性动力学空化}

本节有双重目的：它表明，无界弱解对双曲平衡律方程组可能展现出与本书前面各章所遇不同的另一类非唯一性；同时还表明，归因于真空存在的病态并非气体动力学所独有，它可推广到固体的弹性动力学。

我们考虑无体积力 $b=0$ 时等熵热弹性方程组\eqref{align:ch03-等熵热弹性系统}的 Cauchy 问题，初值为
\begin{align}
F(x,0)=\lambda I,\qquad v(x,0)=0,\qquad x\in\mathbb{R}^3,
\label{align:ch18-1881}
\end{align}
其中 $\lambda$ 为正常数。它有定常解 $F(x,t)=\lambda I$、$v(x,t)=0$（$x\in\mathbb{R}^3$，$t\geqslant0$），相伴的弹性体放置为 $\chi(x,t)=\lambda x$。此外，\cref{theorem:ch05-involution-stability}断言：当 Piola--Kirchhoff 应力 $S$ 经由\cref{theorem:ch05-involution-stability}由秩一凸内能函数 $\varepsilon$ 导出时，上述解在满足熵不等式\cref{theorem:ch05-involution-stability}、局部振荡很小的弱解类中是唯一的。然而，我们此处将看到，若去掉弱解必须是 $L^\infty$ 函数的限制，就存在满足熵不等式的额外解。

我们考虑弹性介质为各向同性固体的情形。Cauchy 问题的新解由如下形式的自相似运动诱导：
\begin{align}
\chi(x,t)=\frac{\phi(s)}{s}\,x,\qquad s=\frac{|x|}{t},
\label{align:ch18-1882}
\end{align}
其中 $\phi$ 是 $[0,\infty)$ 上的 Lipschitz 函数，在某 $\sigma>0$ 对应的区间 $[0,\sigma)$ 与 $(\sigma,\infty)$ 上光滑，且
\begin{align}
\phi(s)>0,\qquad \dot\phi(s)>0,\qquad \ddot\phi(s)>0,\qquad 0<s<\sigma,
\label{align:ch18-1883}
\end{align}
\begin{align}
\dot\phi(\sigma-)=\mu<\lambda,
\label{align:ch18-1884}
\end{align}
\begin{align}
\phi(s)=\lambda s,\qquad \sigma\leqslant s<\infty.
\label{align:ch18-1885}
\end{align}

于是，在 $\chi$ 的 Lagrangian 表示中，物体在区域 $|x|>\sigma t$ 内处于静止，该区域以一道半径为 $\sigma t$、随时间线性增长的前驱球面激波为界。激波后方、区域 $0<|x|<\sigma t$ 内，材料经历光滑变形。然而，沿直线 $x=0$ 有一个奇点，类似于\cref{section:ch09-存在性与唯一性的失效}中遇到的 delta 激波。在 $\chi$ 的 Euler 表示中，该奇线代表一个球形空腔，它在 $t=0$ 时于原点张开，并随其半径 $\phi(0)t$ 线性增长而膨胀。我们要求空腔自发出现并自平衡，使物体的自由界面无应力：Cauchy 应力的径向分量在其上为零。

由运动\eqref{align:ch18-1882}生成的变形梯度场与速度场为
\begin{align}
F=\frac{\phi(s)}{s}I-\left[\frac{\phi(s)}{s}-\dot\phi(s)\right]\frac{1}{|x|^2}xx^\top,
\label{align:ch18-1886}
\end{align}
\begin{align}
v=\left[\frac{\phi(s)}{s}-\dot\phi(s)\right]\frac1t x.
\label{align:ch18-1887}
\end{align}
目标是构造 $\phi$，使\eqref{align:ch18-1886}、\eqref{align:ch18-1887} 提供\eqref{align:ch03-等熵热弹性系统}、\eqref{align:ch18-1881}的解。

首先注意，$(F,v)$ 在分布意义下满足\eqref{align:ch03-等熵热弹性系统}$_1$。事实上，\eqref{align:ch18-1881}$_1$ 在 $0<|x|<\sigma t$ 与 $\sigma t<|x|<\infty$ 上逐点成立；此外，跨曲面 $|x|=\sigma t$ 的跳跃条件\eqref{align:ch03-等熵热弹性RH}$_1$ 作为 $\chi$ 连续性的直接推论成立。最后，$1/s$ 一项在 $x=0$ 处奇异，但不造成任何问题，因为它局部可积。

现在转向表达动量平衡的方程\eqref{align:ch03-等熵热弹性系统}$_2$。如\cref{section:ch02-热弹性}中讨论，任一各向同性热弹性固体在常熵下的内能 $\varepsilon$ 是右伸长张量 $U$（由\eqref{align:ch02-2-1-7}定义）特征值的对称函数
\begin{align}
\varepsilon=\Phi(\lambda_1,\lambda_2,\lambda_3).
\label{align:ch18-1888}
\end{align}
此处 $F$ 正定对称，故 $U=F$，且
\begin{align}
\lambda_1=\dot\phi(s),\qquad \lambda_2=\frac{\phi(s)}{s},\qquad \lambda_3=\frac{\phi(s)}{s}.
\label{align:ch18-1889}
\end{align}
为方便计，我们用记号 $\Phi_i,\Phi_{ij},\ldots$ 表示函数 $\partial\Phi/\partial\lambda_i$、$\partial^2\Phi/\partial\lambda_i\partial\lambda_j$ 等。进一步，以 $\hat\Phi,\hat\Phi_i,\hat\Phi_{ij}$ 表示将 $\Phi,\Phi_i,\Phi_{ij}$ 与\eqref{align:ch18-1889}复合后所得的 $s$ 的函数。

由\eqref{align:ch03-等熵热弹性本构}与\eqref{align:ch18-1888}，经冗长计算得
\begin{align}
S=\hat\Phi_1(s)\left[\frac{1}{|x|^2}xx^\top\right]+\hat\Phi_2(s)\left[I-\frac{1}{|x|^2}xx^\top\right].
\label{align:ch18-18810}
\end{align}
于是，只要\eqref{align:ch03-等熵热弹性系统}$_2$ 在 $0<|x|<\sigma t$ 上逐点成立、且跨曲面 $|x|=\sigma t$ 的跳跃条件\eqref{align:ch03-等熵热弹性RH}$_2$ 成立，它就在分布意义下成立。借助\eqref{align:ch18-1887}与\eqref{align:ch18-18810}，这些条件在此退化为
\begin{align}
\frac12\left[\hat\Phi_{11}(s)-s^2\right]s\ddot\phi(s)=\hat\Phi_2(s)-\hat\Phi_1(s)
+\left[\frac{\phi(s)}{s}-\dot\phi(s)\right]\hat\Phi_{12}(s),
\label{align:ch18-18811}
\end{align}
对 $0<s<\sigma$ 成立；以及
\begin{align}
\sigma^2=\frac{\Phi_1(\lambda,\lambda,\lambda)-\Phi_1(\mu,\lambda,\lambda)}{\lambda-\mu}.
\label{align:ch18-18812}
\end{align}

我们将在假设
\begin{align}
\Phi(\lambda_1,\lambda_2,\lambda_3)=\sum_{i=1}^3 g(\lambda_i)+h(\lambda_1\lambda_2\lambda_3),
\label{align:ch18-18813}
\end{align}
下构造满足上述规格的 $\phi$，其中 $g$、$h$ 是 $(0,\infty)$ 上的光滑函数，满足
\begin{align}
g''(p)>0,\qquad g'''(p)<0,\qquad 0<p<\infty,
\label{align:ch18-18814}
\end{align}
\begin{align}
h''(b)>0,\qquad h'''(b)<0,\qquad 0<b<\infty,
\label{align:ch18-18815}
\end{align}
\begin{align}
\lim_{p\to\infty}\frac{g'(p)}{p}=\gamma\geqslant0,
\label{align:ch18-18816}
\end{align}
\begin{align}
h'(b)\to-\infty\quad(b\to0+),\qquad h'(b)\to\infty\quad(b\to\infty),
\label{align:ch18-18817}
\end{align}
由此得
\begin{align}
\Phi_{11}(p,q,q)>0,\qquad \Phi_{111}(p,q,q)<0,\qquad 0<p<\infty,\ 0<q<\infty.
\label{align:ch18-18818}
\end{align}
上述条件在物理上是合理的，它们刻画了具有材料软化、多凸应变能函数的各向同性弹性固体。

拉伸参数 $\lambda$ 不事先给定，而在构造过程中确定。

我们寻求常微分方程\eqref{align:ch18-18811}在初始条件
\begin{align}
\phi(0)=\phi_0,\qquad \lim_{s\to0+}\dot\phi(s)\left[\frac{\phi(s)}{s}\right]^2=a
\label{align:ch18-18819}
\end{align}
下的解，其中 $\phi_0$ 是自由正参数，而 $a$ 这样固定：使 Cauchy 应力的径向分量在自由边界上为零。将\eqref{align:ch02-2-3-6}与\eqref{align:ch18-18810}联立计算可得：当 $h'(a)=0$ 时，上述要求满足。

初等但艰深的分析（载于\cref{section:ch18-注记}所引文献）表明，\eqref{align:ch18-18811}、\eqref{align:ch18-18819} 在最大区间 $[0,\bar s)$ 上有唯一解 $\phi$，满足
\begin{align}
\phi(s)>0,\quad \dot\phi(s)>0,\quad \ddot\phi(s)>0,\quad \frac{\phi(s)}{s}-\dot\phi(s)>0,\qquad 0<s<\bar s.
\label{align:ch18-18820}
\end{align}
此外，区间 $(0,\bar s)$ 中存在唯一的 $\sigma$，使\eqref{align:ch18-18812}以 $\lambda=\phi(\sigma)/\sigma$、$\mu=\dot\phi(\sigma)$ 成立。这就完成了具有上述规格的运动 $\chi$ 的构造。

分析还表明：变动参数 $\phi_0$，可使拉伸参数 $\lambda$ 取到临界值 $\lambda_{cr}$ 以上的任意值。

现在转向所构造解的可容性。由\eqref{align:ch18-18818}，且因 $\mu<\lambda$，
\begin{align}
\Phi_{11}(\lambda,\lambda,\lambda)<\sigma^2<\Phi_{11}(\mu,\lambda,\lambda),
\label{align:ch18-18821}
\end{align}
它验证激波满足 Lax 的 $E$-条件。

接着检验激波熵可容性条件\eqref{align:ch03-等熵熵不等式}的等价形式\eqref{align:ch03-等熵RH熵条件}。利用\eqref{align:ch18-1886}、\eqref{align:ch18-1888}与\eqref{align:ch18-18810}直接计算得
\begin{align}
\begin{aligned}
[\![\varepsilon]\!]-\operatorname{tr}\left(\frac12(S_++S_-)^\top[\![F]\!]\right)
&=\Phi(\lambda,\lambda,\lambda)-\Phi(\mu,\lambda,\lambda)\\
&\quad-\frac12(\lambda-\mu)[\Phi_1(\lambda,\lambda,\lambda)+\Phi_1(\mu,\lambda,\lambda)].
\end{aligned}
\label{align:ch18-18822}
\end{align}
上式右端是 $\mu$ 的函数，在 $\mu=\lambda$ 处与其一阶导数同时为零，而其二阶导数 $-\frac12(\lambda-\mu)\Phi_{111}(\mu,\lambda,\lambda)$ 由\eqref{align:ch18-18818}知非负。于是结论：激波也满足熵可容性条件\eqref{align:ch03-等熵RH熵条件}。

既然定常解与空化解都满足 Lax 的 $E$-条件与熵条件，Cauchy 问题\eqref{align:ch03-等熵热弹性系统}、\eqref{align:ch18-1881}因这一新解而提出的非唯一性问题，就不能靠诉诸标准可容性准则来解决。于是我们转而比较这两个解的熵产生速率，其精神与\cref{section:ch09-熵率容许性准则}中的讨论一致。

新解 $(F,v)$ 相对于定常解 $(\lambda I,0)$ 的相对熵为
\begin{align}
H=\varepsilon(F)+\frac12|v|^2-\varepsilon(\lambda I)-\operatorname{tr}\left(S(\lambda I)(F-\lambda I)^\top\right).
\label{align:ch18-18823}
\end{align}
因而相对熵产生速率为
\begin{align}
\dot{\mathscr H}(t)=\int_{|x|<\sigma t}\partial_t H(x,t)\,\mathrm{d}x+\sigma\int_{|x|=\sigma t}H(x,t)\,\mathrm{d}a.
\label{align:ch18-18824}
\end{align}
直接计算后，
\begin{align}
\begin{aligned}
\dot{\mathscr H}(t)=-4\pi\sigma^3t^2\bigl\{&\Phi(\lambda,\lambda,\lambda)-\Phi(\mu,\lambda,\lambda)\\
&-\frac12(\lambda-\mu)[\Phi_1(\lambda,\lambda,\lambda)+\Phi_1(\mu,\lambda,\lambda)]\bigr\},
\end{aligned}
\label{align:ch18-18825}
\end{align}
它为负（对照\eqref{align:ch18-18822}）。故熵产生准则更偏爱空化解而非定常解。当然，这一论证并不能挑出唯一解，因为人们可以生成无穷多个解，其中空腔在物体不同点同时张开。因而可构造出熵产生速率任意小的解。

\end{document}
